% Part: first-order-logic
% Chapter: completeness
% Section: lindenbaums-lemma

\documentclass[../../../include/open-logic-section]{subfiles}

\begin{document}

\iftag{FOL}
      {\olfileid{fol}{com}{lin}}
      {\olfileid{pl}{com}{lin}}

\olsection{Lemma van Lindenbaum}

\begin{explain}
We bewijzen nu een lemma waaruit blijkt dat elke consistente verzameling
!!{sentence}s opgenomen is in een verzameling zinnen die niet alleen
consistent, maar ook !!{complete} is. In het bewijs voegen we telkens één
!!{sentence} toe en garanderen we bij elke stap dat de verzameling
consistent blijft. We doen dit zo dat voor elke $!A$ in een zeker stadium
ofwel $!A$ ofwel $\lnot !A$ wordt toegevoegd. De vereniging van alle
stadia in deze constructie bevat dan $!A$ of haar ontkenning~$\lnot !A$
en is dus volledig. Zij is ook consistent, omdat we er bij elke stap voor
zorgen dat we geen inconsistentie invoeren.
\end{explain}

\begin{lem}[Lemma van Lindenbaum]
\ollabel{lem:lindenbaum} Elke consistente verzameling~$\Gamma$ in een
taal~$\Lang{L}$ kan worden uitgebreid tot !!a{complete} en consistente
verzameling~$\Gamma^*$.
\end{lem}

\begin{proof}
Laat $\Gamma$ consistent zijn. Laat $!A_0$, $!A_1$, \dots{} een
opsomming zijn van alle !!{sentence}s van~$\Lang L$. Definieer
$\Gamma_0 = \Gamma$ en
\[
\Gamma_{n+1} =
\begin{cases}
\Gamma_n \cup \{ !A_n \} & \textrm{als $\Gamma_n \cup \{!A_n\}$
  consistent is;} \\
\Gamma_n \cup \{ \lnot !A_n \} & \textrm{anders.}
\end{cases}
\]
Stel $\Gamma^* = \bigcup_{n \geq 0} \Gamma_n$.

Elke $\Gamma_n$ is consistent: $\Gamma_0$ is per definitie consistent.
Als $\Gamma_{n+1} = \Gamma_n \cup \{!A_n\}$, dan is dit omdat de
laatstgenoemde verzameling consistent is. Als zij dat niet is, dan
$\Gamma_{n+1} = \Gamma_n \cup \{\lnot !A_n\}$. We moeten verifiëren dat
$\Gamma_n \cup \{\lnot !A_n\}$ consistent is. Veronderstel dat dit niet
zo is. Dan zijn \emph{zowel} $\Gamma_n \cup \{!A_n\}$ als
$\Gamma_n \cup \{\lnot !A_n\}$ inconsistent. Volgens
\iftag{FOL}{%
  \tagrefs{prfAX/{fol:axd:prv:prop:provability-exhaustive},
    prfSC/{fol:seq:prv:prop:provability-exhaustive},
    prfND/{fol:ntd:prv:prop:provability-exhaustive},
    prfTab/{fol:tab:prv:prop:provability-exhaustive}}}{%
  \tagrefs{prfAX/{pl:axd:prv:prop:provability-exhaustive},
    prfSC/{pl:seq:prv:prop:provability-exhaustive},
    prfND/{pl:ntd:prv:prop:provability-exhaustive},
    prfTab/{pl:tab:prv:prop:provability-exhaustive}}}
zou $\Gamma_n$ dan inconsistent zijn, in tegenspraak met de
inductiehypothese.

Voor elke~$n$ en elke $i < n$ geldt $\Gamma_i \subseteq \Gamma_n$. Dit
volgt door een eenvoudige inductie naar~$n$. Voor $n=0$ bestaan er geen
$i < 0$, zodat de bewering automatisch geldt. Veronderstel voor de
inductiestap dat zij geldt voor~$n$. We tonen aan dat uit $i < n+1$
volgt dat $\Gamma_i \subseteq \Gamma_{n+1}$. Volgens de constructie is
$\Gamma_{n+1} = \Gamma_n \cup \{!A_n\}$ of $= \Gamma_n \cup \{\lnot
!A_n\}$. Dus $\Gamma_n \subseteq \Gamma_{n+1}$. Als $i < n+1$, dan
geldt volgens de inductiehypothese $\Gamma_i \subseteq \Gamma_n$ (als
$i < n$), of geldt het triviale feit $\Gamma_n \subseteq \Gamma_n$ (als
$i = n$). We verkrijgen dan $\Gamma_i \subseteq \Gamma_{n+1}$ door de
transitiviteit van~$\subseteq$.

Hieruit volgt dat $\Gamma^*$ consistent is. Dit kan als volgt worden
ingezien. Laat $\Gamma' \subseteq \Gamma^*$ eindig zijn. Elk
$!B \in \Gamma'$ behoort dan ook tot~$\Gamma_i$ voor een zekere~$i$.
Laat $n$ de grootste van deze indices zijn. Omdat
$\Gamma_i \subseteq \Gamma_n$ als $i \le n$, geldt voor elk
$!B \in \Gamma'$ ook $\in \Gamma_n$; met andere woorden,
$\Gamma' \subseteq \Gamma_n$. Bovendien is $\Gamma_n$ consistent. Elke
eindige deelverzameling $\Gamma' \subseteq \Gamma^*$ is dus consistent.
Volgens \iftag{FOL}{%
  \tagrefs{prfAX/{fol:axd:ptn:prop:proves-compact},
    prfSC/{fol:seq:ptn:prop:proves-compact},
    prfND/{fol:ntd:ptn:prop:proves-compact},
    prfTab/{fol:tab:ptn:prop:proves-compact}}}{%
  \tagrefs{prfAX/{pl:axd:ptn:prop:proves-compact},
    prfSC/{pl:seq:ptn:prop:proves-compact},
    prfND/{pl:ntd:ptn:prop:proves-compact},
    prfTab/{pl:tab:ptn:prop:proves-compact}}} is $\Gamma^*$ consistent.

Elke !!{sentence} van $\Frm[L]$ komt voor in de opsomming waarmee
~$\Gamma^*$ is gedefinieerd. Als $!A_n \notin \Gamma^*$, dan komt dit
doordat $\Gamma_n \cup \{!A_n\}$ inconsistent was. Maar dan geldt
$\lnot !A_n \in \Gamma^*$, zodat $\Gamma^*$ !!{complete} is.
\end{proof}

\end{document}
